\chapter{Ex\'amenes} 

% ==============================================================================
% 14 DE ENERO DE 2026
% ==============================================================================
\section{14 de enero de 2026}

\subsection*{TEOR\'IA}

\begin{enumerate}
    \item Demostrar que si un espacio topológico es conexo y localmente arcoconexo, entonces es arcoconexo. (Demostrar los resultados sobre espacios conexos y arcoconexos que se utilicen).
    \hfill \textbf{(2.5 puntos)}

    \item Si $X$ es un espacio topológico e $Y$ un subespacio, ¿es cierto que los puntos del borde de $Y$ son los puntos del cierre de $Y$ que no pertenecen al interior de $Y$? Razónese la respuesta, en caso afirmativo, demuéstrese y en caso negativo, modificar el enunciado para caracterizar los puntos del borde, demostrándolo.
    \hfill \textbf{(1 punto)}

    \item Si $X$ es un espacio topológico, ¿es cierto que si $X$ es compacto por límites entonces es compacto por sucesiones? Razónese la respuesta, en caso afirmativo, demuéstrese y en caso negativo, añadir las hipótesis necesarias al enunciado anterior para que sea cierto, demostrándolo.
    \hfill \textbf{(1.5 puntos)}
\end{enumerate}

\subsection*{PROBLEMAS}

\begin{enumerate}
    \item Sean $(X, d)$ y $(X', d')$ espacios seudométricos y $\mathcal{T}_d$, $\mathcal{T}_{d'}$ las topologías asociadas a las seudométricas. Sea $f : X \to X'$ una isometría. Demostrar que:
    \begin{enumerate}
        \item La topología $\mathcal{T}_d$ coincide con la topología inicial de la aplicación $f : X \to (X', \mathcal{T}_{d'})$.
        \hfill \textbf{(0.5 puntos)}
        
        \item Si $f$ es epiyectiva, entonces la topología $\mathcal{T}_{d'}$ coincide con la topología final de la aplicación $f : (X, \mathcal{T}_d) \to X'$. ¿Es cierto el enunciado anterior si $f$ no es epiyectiva?
        \hfill \textbf{(1 punto)}
    \end{enumerate}

    \item Sean $X, Y$ espacios topológicos conexos y $Z$ otro espacio topológico. Sea $f : X \times Y \to Z$ una aplicación de conjuntos, tal que para cada $x_0 \in X$ y para cada $y_0 \in Y$, las aplicaciones
    \[
        f_{x_0} : Y \to Z, \quad y \mapsto f(x_0, y)
    \]
    \[
        f_{y_0} : X \to Z, \quad x \mapsto f(x, y_0)
    \]
    son aplicaciones continuas. Demostrar que $\operatorname{Im} f = f(X \times Y) \subset Z$ es un subespacio conexo.
    \hfill \textbf{(1.5 puntos)}

    \item Sea $X$ un espacio topológico regular (no necesariamente $T_2$) e $Y \subset X$ un subespacio compacto. Demostrar que para todo abierto $U \subset X$ con $Y \subset U$, existe un abierto $V \subset X$ tal que $Y \subset V \subset \overline{V} \subset U$.
    \hfill \textbf{(2 puntos)}
\end{enumerate}

\newpage

% ==============================================================================
% 19 DE NOVIEMBRE DE 2025
% ==============================================================================
\section{19 de noviembre de 2025}

\begin{enumerate}
    \item 
    \begin{enumerate}
        \item Enunciar con detalle y precisión todos los axiomas de separación.
        \hfill \textbf{(2 puntos)}
        \item Dar y demostrar un enunciado equivalente no trivial para cada axioma.
        \hfill \textbf{(3 puntos)}
    \end{enumerate}

    \item Sea $(X, d)$ un espacio seudométrico y $\mathcal{T}_d$ la topología inducida por la seudométrica.
    \begin{enumerate}
        \item Demostrar que la aplicación distancia $d : (X, \mathcal{T}_d) \times (X, \mathcal{T}_d) \to \mathbb{R}$ es continua.
        \hfill \textbf{(1 punto)}
        \item Demostrar que la topología $\mathcal{T}_d$ coincide con la topología inicial de las aplicaciones $\{d_x : X \to \mathbb{R}\}_{x \in X}$ definidas por $d_x(x') := d(x, x')$ para todo $x' \in X$.
        \hfill \textbf{(2 puntos)}
        \item Sea $\mathcal{T}$ otra topología en $X$. Demostrar que la topología $\mathcal{T}_d$ es menos fina que la topología $\mathcal{T}$ si y solo si la aplicación distancia $d : (X, \mathcal{T}) \times (X, \mathcal{T}) \to \mathbb{R}$ es continua.
        \hfill \textbf{(2 puntos)}
    \end{enumerate}
\end{enumerate}

\newpage

% ==============================================================================
% 16 DE ENERO DE 2025
% ==============================================================================
\section{16 de enero de 2025}

\subsection*{TEOR\'IA}

\begin{enumerate}
    \item Define con precisión espacio topológico compacto, espacio topológico compacto por límite y espacio topológico compacto por sucesiones. Demuestra que:
    \hfill \textbf{(2.5 puntos)}
    \begin{enumerate}
        \item Todo espacio compacto es compacto por límite.
        \item Todo espacio metrizable compacto por límite es compacto por sucesiones.
    \end{enumerate}

    \item Define con precisión espacio topológico conexo, espacio topológico arcoconexo, espacio topológico localmente conexo y espacio topológico localmente arcoconexo. Demuestra que:
    \hfill \textbf{(2.5 puntos)}
    \begin{enumerate}
        \item Las componentes conexas de un espacio localmente conexo son abiertas y cerradas.
        \item Las componentes arcoconexas de un espacio localmente arcoconexo son abiertas y cerradas.
    \end{enumerate}
\end{enumerate}

\subsection*{PROBLEMAS}

\begin{enumerate}
    \item Sea $X$ un espacio topológico y $D \subset X$ un subespacio denso. Demostrar que:
    \begin{enumerate}
        \item Para todo abierto $U \subset X$ se verifica que $\overline{U} = \overline{U \cap D}$.
        \hfill \textbf{(0.75 puntos)}
        \item Si $V \subset D$ es un abierto de $D$, existe un abierto $U$ de $X$ tal que el cierre de $V$ en $X$ y el cierre de $U$ coinciden.
        \hfill \textbf{(0.25 puntos)}
        \item Deducir que si $X$ es separado y $D$ es localmente compacto, entonces $D$ es abierto en $X$.
        \hfill \textbf{(0.75 puntos)}
    \end{enumerate}

    \item Estudiar si son conexos los siguientes subconjuntos de $\mathbb{R}^2$ con la topología usual:
    \[
        \mathbb{Z} \times \mathbb{Z} \quad \text{y} \quad (\mathbb{Z} \times \mathbb{R}) \cup (\mathbb{R} \times \mathbb{Z})
    \]
    En caso negativo, calcular sus componentes conexas.
    \hfill \textbf{(1.5 puntos)}

    \item Sea $(X, d)$ un espacio métrico. Demostrar:
    \begin{enumerate}
        \item Para todo subconjunto $A \subset X$, $x \in \overline{A}$ si y solo si $d(x, A) = 0$.
        \hfill \textbf{(0.5 puntos)}
        \item El espacio métrico $X$ es completamente normal, es decir, para cualesquiera subconjuntos $A, B \subset X$ tales que $\overline{A} \cap B = A \cap \overline{B} = \emptyset$ existen abiertos $U, V \subset X$ tales que $A \subset U$, $B \subset V$ y $U \cap V = \emptyset$.
        \hfill \textbf{(1.25 puntos)}
    \end{enumerate}
\end{enumerate}

\newpage

% ==============================================================================
% 6 DE NOVIEMBRE DE 2024
% ==============================================================================
\section{6 de noviembre de 2024}

\begin{enumerate}
    \item Sea $X$ un conjunto y $\mathcal{B} \subset \mathcal{P}(X)$ una familia de subconjuntos. Enunciar y demostrar las propiedades necesarias y suficientes que ha de verificar $\mathcal{B}$ para que sea la base de una topología por cerrados en $X$.
    \hfill \textbf{(3.5 puntos)}

    \item Sea $X$ un espacio topológico y $A \subset X$ un subconjunto. Demostrar que si $A$ es abierto o cerrado entonces el interior de la frontera de $A$ es el vacío. ¿Es cierto el recíproco?
    \hfill \textbf{(2.5 puntos)}

    \item Sea $X$ un espacio topológico y $A \subset X$ un subconjunto. Demostrar que $A$ es abierto si y solo si verifica que para todo subconjunto $B \subset X$ tal que $A \cap B = \emptyset$ entonces $A \cap \overline{B} = \emptyset$.
    \hfill \textbf{(2 puntos)}

    \item Sean $f : X \to Y$ y $g : Y \to Z$ aplicaciones entre espacios topológicos. Teniendo en cuenta que una aplicación abierta o cerrada no es necesariamente continua, demostrar:
    \begin{enumerate}
        \item Si $g \circ f$ es abierta y $f$ es continua y epiyectiva, entonces $g$ es abierta.
        \hfill \textbf{(1 punto)}
        \item Si $g \circ f$ es cerrada y $g$ es continua e inyectiva, entonces $f$ es cerrada.
        \hfill \textbf{(1 punto)}
    \end{enumerate}
\end{enumerate}

\newpage

% ==============================================================================
% 30 DE ENERO DE 2024
% ==============================================================================
\section{30 de enero de 2024}

\subsection*{TEORÍA}

\begin{enumerate}
    \item Sea $(X, d)$ un espacio métrico.
    \begin{enumerate}
        \item Sea $A \subset X$ un subconjunto. Demostrar que el cierre de $A$ es $\overline{A} = \{x \in X : d(x, A) = 0\}$.
        \hfill \textbf{(1 punto)}
        \item Si $Y$ es un espacio topológico metrizable, demostrar que una aplicación $f : X \to Y$ es continua si y solo si para toda sucesión $\{x_n\} \subset X$ que converja a un punto $x \in X$ se verifica que la sucesión $\{f(x_n)\} \subset Y$ converge al punto $f(x) \in Y$. ¿Es necesaria la hipótesis de que $Y$ sea metrizable? Razónese la respuesta.
        \hfill \textbf{(1 punto)}
    \end{enumerate}

    \item Sea $X$ un espacio localmente compacto y separado. Construir como espacio topológico su compactificación de Alexandroff y enunciar y demostrar sus propiedades.
    \hfill \textbf{(3 puntos)}
\end{enumerate}

\subsection*{PROBLEMAS}

\begin{enumerate}
    \item Sea $X$ un espacio topológico separado y $U \subset X$ un abierto no vacío.
    \begin{enumerate}
        \item Demostrar que para todo punto $x \notin U$, existe un abierto $V \subset X$ tal que $V \subset U$ y $x \notin \overline{V}$.
        \hfill \textbf{(1 punto)}
        \item Si $x \in U$, ¿es cierta la afirmación del apartado anterior? En caso afirmativo, demuéstrese, y en caso negativo, dar un contraejemplo y dar una condición topológica sobre el espacio $X$ o sobre el punto $x$ considerado para que sea cierta la afirmación cuando $x \in U$, y en este caso, demuéstrese.
        \hfill \textbf{(1 punto)}
    \end{enumerate}

    \item Sea $\{X_i\}_{i \in I}$ una familia arbitraria de espacios topológicos. Demostrar que si cada espacio $X_i$ es $T_1$ para todo $i \in I$, entonces el producto directo $\prod_{i \in I} X_i$ es un espacio topológico $T_1$.
    \hfill \textbf{(1.5 puntos)}

    \item Sean $X$ un espacio topológico y $\{Y_i\}_{i \in I}$ una familia de subespacios conexos de $X$. Probar que si existe un $k \in I$ tal que $\overline{Y_k} \cap Y_i \neq \emptyset$ para todo $i \in I$, entonces $\bigcup_{i \in I} Y_i$ es conexo.
    \hfill \textbf{(1.5 puntos)}
\end{enumerate}

\newpage

% ==============================================================================
% 11 DE ENERO DE 2024
% ==============================================================================
\section{11 de enero de 2024}

\subsection*{TEORÍA}

\begin{enumerate}
    \item Sea $(X, \mathcal{T})$ un espacio topológico siendo $\mathcal{T}$ su topología por abiertos. Definir los entornos de un punto de $X$. Enunciar y demostrar las condiciones necesarias y suficientes para que unas familias de subconjuntos $\{\mathcal{E}'(x) \subset \mathcal{P}(X)\}_{x \in X}$ sean las familias de entornos de una topología en $X$. Demostrar que esta topología es única.
    \hfill \textbf{(2 puntos)}

    \item Dar con precisión las definiciones de espacio topológico compacto, localmente compacto, compacto por límites y compacto por sucesiones. Enunciar con claridad las relaciones que existen entre ellos.
    \hfill \textbf{(2 puntos)}

    \item Sea $(X, d)$ un espacio métrico completo. Demostrar que un subconjunto $Y \subset X$ es completo si y sólo si $Y$ es cerrado por la topología inducida por $d$.
    \hfill \textbf{(1 punto)}
\end{enumerate}

\subsection*{PROBLEMAS}

\begin{enumerate}
    \item Sea $X$ un espacio topológico y $R$ la relación de equivalencia definida por
    \[
        x Ry \quad \text{si} \quad \overline{\{x\}} = \overline{\{y\}} \quad \text{con } x, y \in X.
    \]
    Sea $\pi : X \to X/R$ la aplicación de paso al cociente y $X/R$ espacio topológico con la topología cociente. Demostrar que:
    \begin{enumerate}
        \item Todo cerrado de $X$ es saturado para la relación $R$.
        \hfill \textbf{(0.75 puntos)}
        \item Todo abierto de $X$ es saturado para la relación $R$.
        \hfill \textbf{(0.75 puntos)}
        \item La aplicación $\pi$ es abierta y cerrada. ¿Es homeomorfismo? Razónese la respuesta.
        \hfill \textbf{(0.5 puntos)}
    \end{enumerate}

    \item Sea $X$ un espacio topológico compacto y $A \subseteq X$ un subconjunto infinito. Entonces el subconjunto $A'$ de los puntos de acumulación de $A$ es no vacío.
    \hfill \textbf{(1.5 puntos)}

    \item Sea $X$ un espacio métrico en el que toda bola cerrada es compacta. Demostrar que $X$ es completo. Deducir que $\mathbb{R}^n$ es completo con la métrica usual.
    \hfill \textbf{(1.5 puntos)}
\end{enumerate}

\newpage

% ==============================================================================
% 15 DE NOVIEMBRE DE 2023
% ==============================================================================
\section{15 de noviembre de 2023}

\begin{enumerate}
    \item Demostrar que $\mathbb{R}^n$ con la topología producto directo es metrizable.
    \hfill \textbf{(3.5 puntos)}

    \item ¿Es cierto que si $f : X \to Y$ es una aplicación entre espacios topológicos y $X$ tiene la topología inicial de $f$, entonces $f$ es una aplicación continua abierta? En caso afirmativo, probarlo, y en caso negativo, dar un contraejemplo y dar una condición sobre la aplicación o sobre los espacios topológicos para que sea cierta la afirmación anterior.
    \hfill \textbf{(2.5 puntos)}

    \item Sea $X$ el intervalo abierto $(0, 1) \subset \mathbb{R}$ y
    \[
        \mathcal{T} = \left\{ U_n = \left(0, 1 - \frac{1}{n}\right) \right\}_{n=2, 3, \dots} \cup \{\emptyset, X\}.
    \]
    Consideramos los subconjuntos $Y_1 = \left(\frac{2}{3}, \frac{5}{6}\right]$ y $Y_2 = \left(0, \frac{5}{7}\right)$ de $X$.
    \begin{enumerate}
        \item Probar que $\mathcal{T}$ es una topología por abiertos en $X$.
        \item Calcular los cierres, interiores y fronteras de los subespacios $Y_1, Y_2$.
        \item Calcular el exterior de $Y_1$ y el borde de $Y_2$.
        \item Calcular los subespacios conexos de $X$ y sus componentes conexas.
        \item Estudiar los axiomas de separación en $(X, \mathcal{T})$, es decir, si es $T_0, T_1, T_2, T_3, T_4$. ¿El subespacio $\{\frac{1}{4}\}$ es $T_0$? ¿Y el subespacio $\{\frac{1}{4}, \frac{2}{3}\}$ es $T_0$?
        \item Comparar la topología $\mathcal{T}_d$ métrica usual de $X$ y la topología $\mathcal{T}$ dada en $X$.
        \item Probar que la aplicación $f : ((1, 2), \mathcal{T}_d) \to (X, \mathcal{T})$ dada por $f(x) = x - 1$ es continua.
    \end{enumerate}
    \hfill \textbf{(4 puntos)}
\end{enumerate}

% ==============================================================================
\section{10 de enero de 2023}
% ==============================================================================

\subsection*{TEORÍA}
\begin{enumerate}
    \item Definir espacio topológico separado. Probar que $X$ es separado si y sólo si la diagonal de $X \times X$ es un cerrado. \hfill (3 puntos)
    \item Sea $(X, d)$ un espacio métrico. Definir la completación $\widehat{X}$ de $X$ y probar que $\widehat{X}$ es un espacio métrico completo (probándose los resultados previos que se requieran). \hfill (5 puntos)
    \item Sea $n > 0$. Probar que el espacio proyectivo real $\mathbb{P}^n(\mathbb{R})$ es compacto. \hfill (2 puntos)
\end{enumerate}

\subsection*{PROBLEMAS}
\begin{enumerate}
    \item Sea $f: X \to Y$ una aplicación entre dos espacios topológicos. Probar que:
    \begin{enumerate}[label=(\alph*)]
        \item $f$ es continua si y sólo si $f^{-1}(A^\circ) \subseteq (f^{-1}(A))^\circ$ para todo $A \subseteq Y$. \hfill (1 punto)
        \item $f$ es cerrada si y sólo si $\overline{f(A)} \subseteq f(\overline{A})$ para todo $A \subseteq X$. \hfill (1 punto)
        \item Si $X$ es compacto y $Y$ es separado ($T_2$), entonces $f$ es un homeomorfismo si y sólo si $f$ es continua y biyectiva. \hfill (1 punto)
        \item Si $X$ y $Y$ tienen la topología cofinita, entonces $f$ es un homeomorfismo si y sólo si $f$ es biyectiva. \hfill (1 punto)
    \end{enumerate}
    
    \item Sea $(\mathbb{R}, \mathcal{T}_{\inf})$ donde $\mathcal{T}_{\inf}$ es la topología generada por la base de abiertos 
    \[
    \mathcal{B} = \{[a, b) \mid a, b \in \mathbb{R}, a < b\}.
    \]
    \begin{enumerate}[label=(\alph*)]
        \item Decidir (de manera justificada) si $(\mathbb{R}, \mathcal{T}_{\inf})$ es conexo o no. Calcula sus componentes conexas. ¿Es totalmente desconexo? \hfill (3 puntos)
        \item Estudiar la compacidad y compacidad local de $(\mathbb{R}, \mathcal{T}_{\inf})$. \hfill (2 puntos)
        \item Sea $\mathcal{T}_u$ la topología usual de $\mathbb{R}$. ¿Es cierto que las únicas aplicaciones $f: (\mathbb{R}, \mathcal{T}_u) \to (\mathbb{R}, \mathcal{T}_{\inf})$ continuas son las constantes? \hfill (1 punto)
    \end{enumerate}
\end{enumerate}

\newpage

% ==============================================================================
\section{30 de enero de 2023}
% ==============================================================================

\subsection*{TEORÍA}
\begin{enumerate}
    \item Definir el grupo fundamental de un espacio topológico $X$ en un punto $p$. Definir cuándo dos espacios topológicos $X, Y$, con puntos prefijados $x_0 \in X, y_0 \in Y$, son homótopos. Demostrar que si $(X, x_0)$ y $(Y, y_0)$ son homótopos, entonces sus grupos fundamentales son isomorfos. \hfill (10 puntos)
    \item Definir espacio topológico compacto. Dados dos espacios topológicos $X, Y$, probar que $X \times Y$ es compacto si y sólo si $X$ y $Y$ son compactos. \hfill (10 puntos)
\end{enumerate}

\subsection*{PROBLEMAS}
\begin{enumerate}
    \item Sea $(X, \mathcal{T})$ un espacio topológico y $A \subseteq X$. Probar que $A^\circ \neq \emptyset$ si y sólo si $A \cap D \neq \emptyset$ para todo conjunto denso $D \subseteq X$. \hfill (10 puntos)
    \item Sea $\mathbb{R}$ con la topología usual y sean $\mathbb{Q}$ y $\mathbb{N}$ con las topologías inducidas.
    \begin{enumerate}[label=(\alph*)]
        \item ¿Es cierto que las únicas aplicaciones $f: \mathbb{R} \to \mathbb{Q}$ continuas son las constantes? \hfill (5 puntos)
        \item Sea $\mathbb{N}$ con la topología inducida por la usual de $\mathbb{R}$. Decidir si la aplicación $f: \mathbb{Q} \to \mathbb{N}$ definida por $f(m/n) = n$, siendo $m/n$ irreducible, es continua o no. \hfill (5 puntos)
    \end{enumerate}
    \item Sean $(X, \mathcal{T})$ un espacio topológico, $\{x_n\}_{n \in \mathbb{N}}$ una sucesión en $X$ y $x \in X$ un límite de dicha sucesión. Probar que $K = \{x_n\}_{n \in \mathbb{N}} \cup \{x\}$ es compacto. \hfill (10 puntos)
\end{enumerate}

\newpage

% ==============================================================================
\section{8 de noviembre de 2022}
% ==============================================================================

\subsection*{TEORÍA (Verdadero o Falso)}
\textit{Contestar verdadero o falso. Cada pregunta acertada vale 1 punto. Cada pregunta mal contestada resta 1 punto. Cada pregunta sin contestar vale 0 puntos.}

\begin{enumerate}
    \item Sea $X$ un espacio topológico finito. Si $X$ es $T_1$, entonces $X$ tiene la topología discreta.
    \item Sean $\mathcal{B}$ y $\mathcal{B}'$ bases para las topologías $\mathcal{T}$ y $\mathcal{T}'$, respectivamente, sobre $X$. Si $\mathcal{T}'$ es más fina que $\mathcal{T}$, entonces para cada $x \in X$ y cada $B \in \mathcal{B}$ con $x \in B$, existe algún $B' \in \mathcal{B}'$ tal que $x \in B' \subseteq B$.
    \item Sea $X$ un espacio topológico y $A \subseteq X$ un subconjunto. Si $A$ no tiene puntos aislados, entonces su cierre, $\overline{A}$, tampoco.
    \item Sean $X$ e $Y$ dos espacios topológicos y $X \times Y$ el espacio topológico producto. Todo cerrado de $X \times Y$ es intersección de cerrados de la forma $A \times B$, con $A$ cerrado de $X$ y $B$ cerrado de $Y$.
    \item Sea $f: X \to Y$ una aplicación continua y epiyectiva entre espacios topológicos. Si la topología de $X$ es la topología inicial de $f$, entonces $Y$ tiene la topología final de $f$.
    \item Sea $X$ un espacio topológico, con una topología $\mathcal{T}$. Sea $d$ una distancia en $X$ y sea $\mathcal{T}_d$ la topología asociada a dicha distancia. Si para todo $x \in X$ y todo $r > 0$ la bola $B(x,r)$ es un abierto de $\mathcal{T}$, entonces $\mathcal{T} = \mathcal{T}_d$.
    \item Sea $X$ un espacio topológico y sea $l: [0, 1] \to X$ un lazo con origen y final en $p \in X$. Si $l$ es homótopo al lazo constante $p$, entonces $l$ y $l^{-1}$ son homótopos.
    \item Sea $Y$ un espacio topológico y $y \in Y$ un punto denso. Para toda aplicación $f: X \to Y$ continua su gráfica $\Gamma_f = \{(x, f(x)) \mid x \in X\}$ es un cerrado de $X \times Y$.
    \item Sea $f: X \to Y$ una aplicación continua. Si $X$ es conexo e $Y$ es discreto, entonces $f$ es constante.
    \item Sean $\mathbb{N}$ y $\mathbb{Q}$, con la topología inducida por la usual de $\mathbb{R}$, entonces no existen aplicaciones biyectivas $\mathbb{Q} \to \mathbb{N}$ continuas.
\end{enumerate}

\subsection*{PROBLEMAS}
\begin{enumerate}
    \item En $\mathbb{R}$ se considera el subconjunto $A = [0, 2) \cup (2, 4) \cup \{6\}$ con la topología $\mathcal{T}_A$ inducida por la topología usual de $\mathbb{R}$.
    \begin{enumerate}[label=(\alph*)]
        \item Decidir si $\{6\}$ es abierto y si es cerrado en $(A, \mathcal{T}_A)$.
        \item Decidir si $(2, 4)$ es abierto y si es cerrado en $(A, \mathcal{T}_A)$.
        \item Calcular los puntos de adherencia de $[0, 2)$ en $(A, \mathcal{T}_A)$.
        \item Decidir si $[0, 1]$ es un entorno de $\{0\}$ en $(A, \mathcal{T}_A)$.
    \end{enumerate}
    \item Sea $f: X \to Y$ una aplicación biyectiva. Probar que $f$ es un homeomorfismo si y sólo si la topología de $Y$ es la topología final de la aplicación $f$.
\end{enumerate}

\newpage

% ==============================================================================
\section{2 de febrero de 2022}
% ==============================================================================

\subsection*{PROBLEMAS}
\begin{enumerate}
    \item Sean $\mathcal{T}$ la topología en $\mathbb{R} \setminus \{0\}$ inducida por la topología usual de $\mathbb{R}$, $\mathcal{T}_{\inf}$ la topología del límite inferior (o de Sorgenfrey) en $\mathbb{R}$ y $f: (\mathbb{R} \setminus \{0\}, \mathcal{T}) \to (\mathbb{R}, \mathcal{T}_{\inf})$ la aplicación definida por $f(x) = 1/x$.
    \begin{enumerate}[label=(\alph*)]
        \item ¿Es $f$ continua? \hfill (8 puntos)
        \item Sea $\beta = \{f^{-1}([a, b)) \text{ con } a < b \text{ números reales}\}$. ¿Es $\beta$ una base para una topología en $\mathbb{R} \setminus \{0\}$? Si $\mathcal{T}_\beta$ es la topología definida por $\beta$ y $g: (\mathbb{R} \setminus \{0\}, \mathcal{T}_\beta) \to (\mathbb{R}, \mathcal{T}_{\inf})$ la aplicación definida por $g(x) = \frac{1}{x}$, deducir que sobre $\text{Im}\,g$ coinciden la topología final de $g$ y la inducida por $\mathcal{T}_{\inf}$. \hfill (8 puntos)
        \item Calcular el cierre y el interior de $A = \{1/n \mid n \in \mathbb{N}\} \subset (\mathbb{R}, \mathcal{T}_{\inf})$. (Indicación: $\mathbb{R}$ es arquimediano, esto es, para todo $z \in \mathbb{R}^+$ existe $n \in \mathbb{N}$ tal que $0 < 1/n < z$). \hfill (8 puntos)
        \item ¿Es $(\mathbb{R} \setminus \{0\}, \mathcal{T}_\beta)$ conexo? Calcular sus componentes conexas. \hfill (8 puntos)
    \end{enumerate}
    
    \item Sea $X$ un espacio topológico compacto e $Y$ un espacio topológico $T_2$.
    \begin{enumerate}[label=(\alph*)]
        \item Probar que para toda aplicación continua y epiyectiva $f: X \to Y$ la topología final de $f$ coincide con la de $Y$. \hfill (8 puntos)
        \item Si $\{K_n\}_{n \in \mathbb{N}}$ es una familia decreciente de subconjuntos compactos de $Y$ (es decir, $K_{i+1} \subseteq K_i$, para todo $i \in \mathbb{N}$), probar que $\displaystyle \bigcap_{n \in \mathbb{N}} K_n \neq \emptyset$. \hfill (8 puntos)
    \end{enumerate}
\end{enumerate}

\newpage

% ==============================================================================
\section{13 de enero de 2022}
% ==============================================================================

\subsection*{TEORÍA}
\begin{enumerate}
    \item Definir espacio compacto y espacio compacto por sucesiones. Si $X$ es un espacio métrico, probar que $X$ es compacto si y sólo si es compacto por sucesiones (pruébense los resultados previos que se requieran). \hfill (14 puntos)
    \item Definir sucesión de Cauchy en un espacio métrico. Definir espacio métrico completo. Dado un espacio métrico $X$, construir su completación y probar que éste es un espacio métrico y completo (pruébense los resultados previos que se requieran). \hfill (14 puntos)
\end{enumerate}

\subsection*{PROBLEMAS}
\begin{enumerate}
    \item Sea $X$ un conjunto, con infinitos elementos, dotado de la topología cofinita 
    \[
    \mathcal{T}_{CF} = \{\emptyset\} \cup \{U \subseteq X \mid X \setminus U \text{ tiene un número finito de elementos}\}.
    \]
    \begin{enumerate}[label=(\alph*)]
        \item Sea $Y \subseteq X$ un subconjunto de $X$ y $y \in Y$. Probar que $Y$ es un entorno de $y$ si y sólo si $X \setminus Y$ tiene un número finito de elementos. \hfill (3 puntos)
        \item Probar que una aplicación $f: (X, \mathcal{T}_{CF}) \to (X, \mathcal{T}_{CF})$ es un homeomorfismo si y sólo si es biyectiva. \hfill (6 puntos)
        \item ¿Es $(X, \mathcal{T}_{CF})$ compacto? ¿Y localmente compacto? Justificar las respuestas. \hfill (6 puntos)
    \end{enumerate}

    \item Sea $(X, \tau)$ un espacio topológico compacto.
    \begin{enumerate}[label=(\alph*)]
        \item ¿Pueden existir subconjuntos $Y \subseteq X$, con infinitos elementos, tal que $\tau_Y$ sea la topología discreta en $Y$? Justificar la respuesta. \hfill (6 puntos)
        \item ¿Pueden existir cerrados $C \subseteq X$, con infinitos elementos, tales que $\tau_C$ sea la topología discreta en $C$? Justificar las respuesta. \hfill (6 puntos)
    \end{enumerate}

    \item Sea $(X, \tau)$ un espacio topológico. Probar que las siguientes afirmaciones son equivalentes:
    \begin{enumerate}[label=(\alph*)]
        \item $X$ es irreducible.
        \item Todo abierto $U \subseteq X$ no vacío es denso.
        \item Todo abierto $U \subseteq X$ no vacío es conexo.
        \item Todo abierto $U \subseteq X$ no vacío es irreducible.
    \end{enumerate} \hfill (15 puntos)
\end{enumerate}

\newpage

% ==============================================================================
\section{10 de noviembre de 2021}
% ==============================================================================

\subsection*{TEORÍA}
\begin{enumerate}
    \item Sea $f: X \to Y$ una aplicación continua y densa (el cierre de la imagen de $f$ es $Y$). Probar que si $X$ es conexo, entonces $Y$ es conexo. Deducir que si $S$ es un subconjunto conexo de un espacio topológico $X$, entonces el cierre de $S$ es conexo.
    \item 
    \begin{enumerate}[label=(\alph*)]
        \item Definir la topología inicial de una aplicación. Definir la topología final de una aplicación.
        \item Sea $f: X \to Y$ una aplicación entre espacios topológicos. Probar que si la topología de $X$ es la inicial de $f$, entonces $g: Z \to X$ es continua si y sólo si $f \circ g: Z \to Y$ lo es. Probar que si la topología de $Y$ es la final de $f$, entonces $g: Y \to Z$ es continua si y sólo si $g \circ f: X \to Z$ lo es.
    \end{enumerate}
\end{enumerate}

\subsection*{PROBLEMAS}
\begin{enumerate} % Inicia en el problema 3
    \item En el conjunto $\mathbb{N} = \{0, 1, 2, \dots\}$ de los números naturales se considera la siguiente familia de subconjuntos:
    \[
    \mathcal{T} = \{\emptyset, \mathbb{N}\} \cup \{\{0, 1, \dots, n\} \mid n \in \mathbb{N}\}.
    \]
    \begin{enumerate}[label=(\alph*)]
        \item Probar que $\mathcal{T}$ es una topología (por abiertos) en $\mathbb{N}$.
        \item Calcula el cierre e interior de $\{0, 1, 2, 3, 7, 10\}$.
        \item Calcula la frontera de $\{4, 13, 23, 37\}$.
        \item Sea $A = \{x_1, \dots, x_r\}$ un conjunto finito de números naturales. Calcular sus puntos de acumulación y puntos aislados.
        \item Sean $f, g: (\mathbb{N}, \mathcal{T}) \to (\mathbb{R}, T_{\inf})$ dos aplicaciones continuas, siendo $T_{\inf}$ la topología del límite inferior (o de Sorgenfrey). Si $f(0) = g(0)$, ¿es cierto que $f = g$? En caso negativo dar un contraejemplo.
    \end{enumerate}
\end{enumerate}